\documentclass[11pt]{article}
\usepackage[margin=1.1in]{geometry}
\usepackage{amsmath,amssymb,amsthm}
\usepackage[T1]{fontenc}
\usepackage{lmodern}
\usepackage{hyperref}
\hypersetup{colorlinks=true,linkcolor=blue,urlcolor=blue,citecolor=blue}
\newtheorem{theorem}{Theorem}
\newtheorem{lemma}[theorem]{Lemma}
\newtheorem{proposition}[theorem]{Proposition}
\theoremstyle{definition}
\newtheorem{definition}[theorem]{Definition}
\setlength{\parskip}{2pt}
\title{A proof of the conjectured linear recurrence and generating function for OEIS A125468}
\author{Adrian Perez Fontelles and Gaspard Moulinier, Independent researchers}
\date{09 September 2026}
\begin{document}
\maketitle
\section{The conjecture}

The Online Encyclopedia of Integer Sequences records \href{https://oeis.org/A125468}{A125468} as

\begin{quote}\itshape Number of base 11 circular n-digit numbers with adjacent digits differing by 9 or less\end{quote}

and carries in its Formula section the lines:

\begin{quote}\ttfamily G.f.: (1 - x\textasciicircum{}2 - 18*x\textasciicircum{}3) / ((1 - x)*(1 - 10*x - 9*x\textasciicircum{}2)).\end{quote}

\begin{quote}\ttfamily a(n) = 11*a(n-1) - a(n-2) - 9*a(n-3) for n>3.\end{quote}

That line carries no separate signature; the entry itself is by R. H. Hardin (Dec 28 2006).

The word {\itshape Empirical} --- equivalently {\itshape Conjecture} --- is the entry's own: the statement is recorded as fitted to the published terms, and no proof is given there. As of revision 17, last modified Aug 12 2026, the entry records no proof and links to none, so the statement is open. This note settles it.

Written out, the claim is

\[\sum_{n \ge 0} a(n)\,x^n \;=\; \frac{1 - x^{2} - 18x^{3}}{1 - 11x + x^{2} + 9x^{3}}.\]

Written out, the claim is

\[a(n) = 11\,a(n-1) - a(n-2) - 9\,a(n-3).\]

stated for $n > 3$.

\section{The object}

The name is read as follows. The reading is not a guess: it is pinned against the entry's own published terms, and against an independent enumeration, before anything is proved (Section 7).

The objects are the words $d_1 d_2 \cdots d_n$ over the digit set $\{0,1,\dots,10\}$ with

\[|d_i - d_{i+1}| \le 9\quad (1 \le i < n), \qquad |d_n - d_1| \le 9.\]

Leading zeros are allowed. This is not an assumption: the entry's own term at $n=1$ is $11$, the number of digits, and not $10$.

The entry's word ``circular'' is what adds the wrap-around pair; the entry's terms at $n = 1$ and $n = 2$ fix that reading as well.

$a(n)$ is the number of such arrays. The entry's offset is 0, so its term of index $n$ is the count for $n$ rows.

\section{The count is a walk count}

Let $G$ be the digraph on the vertex set $\{0,\dots,10\}$ with an edge from $i$ to $j$ exactly when $|i-j| \le 9$, and let $M$ be its adjacency matrix, a $11 \times 11$ symmetric 0--1 band matrix with $M_{ij} = 1$ iff $|i-j| \le 9$.

A circular word $d_1\cdots d_n$ satisfying the condition is exactly a closed walk of length $n$ in $G$, and every closed walk of length $n$ arises from exactly one such word. Hence

\[a(n) \;=\; \operatorname{tr} M^{\,n} \qquad (n \ge 1).\]

\section{The degree bound, derived}

By the Cayley--Hamilton theorem $M$ satisfies its own characteristic polynomial, which has degree $11$. Both $\operatorname{tr} M^n$ and $\mathbf 1^{\mathsf T} M^{n-1} \mathbf 1$ are linear functionals of $M^{n}$, so each satisfies the linear recurrence that polynomial gives. Hence $a$ is C-finite of order at most $S = 11$.

This bound is the size of the matrix the entry's own name determines; it is computed, not assumed. Everything below uses it and nothing else.

\section{The decision procedure}

Let $c_1,\dots,c_D$ be the coefficients of a claimed recurrence $a(n) = \sum_{i=1}^{D} c_i\,a(n-i)$, let $q(t) = t^{D} - \sum_{i=1}^{D} c_i t^{D-i}$ be its characteristic polynomial, and put

\[u_m \;=\; \tilde w^{\mathsf T} L^{\,m-1} q(L)\,\tilde f \qquad (m \ge 1).\]

Expanding the powers of $L$ and using the displayed formula for $a$, one gets $u_m = a(n) - \sum_{i=1}^{D} c_i\,a(n-i)$ with $n = m + D + -1$. So $u_m$ is exactly the residual of the claim at that index, and the recurrence holds at an index $n$ if and only if the corresponding $u_m$ vanishes.

Now $u_m = \tilde w^{\mathsf T} L^{m-1} y$ with $y = q(L)\tilde f$ a fixed vector, so $(u_m)$ satisfies the characteristic polynomial of $L$, of degree $11$: writing that polynomial as $t^{11} - \sum_{i=1}^{11} \lambda_i t^{11-i}$ gives $u_{m+11} = \sum_{i=1}^{11} \lambda_i u_{m+11-i}$. Hence $11$ consecutive vanishing residuals force every later residual to vanish, by induction. The test is therefore finite; and because it locates the {\itshape last} nonvanishing residual, it pins the threshold exactly rather than merely bounding it.

Everything is carried out in exact integer arithmetic on the vector $L^{m-1}y$. The matrix $M$ is never formed.

\section{Results}

\begin{theorem} For A125468, the generating function is exactly $\sum_{n\ge 0} a(n)x^n = \frac{1 - x^{2} - 18x^{3}}{1 - 11x + x^{2} + 9x^{3}}$.\end{theorem}

{\itshape Proof.} Let $N(x)$ and $D(x)$ be the numerator and denominator above, $D(0)=1$. The residual test applied to the recurrence read off $D$ --- namely $a(n) = \sum_i c_i a(n-i)$ with $c_i$ the negated coefficients of $D$ --- gives vanishing residuals throughout, so the power series $F(x)=\sum_{n\ge 0}a(n)x^n$ satisfies $[x^k]\bigl(F(x)D(x)\bigr) = 0$ for every $k > 4$. The remaining coefficients of $F(x)D(x)$, for $k \le 9$, were computed from the walk count and agree with $N(x)$ term by term. Hence $F(x)D(x) = N(x)$ identically, which is the assertion. $\square$

\begin{theorem} For A125468, the recurrence $a(n) = 11\,a(n-1) - a(n-2) - 9\,a(n-3)$ holds for every $n \ge 4$ --- at every index at which it can be stated.\end{theorem}

{\itshape Proof.} The residuals $u_m$ were computed exactly for $m = 1,\dots,38$. Every one of them vanished. After it, more than $S = 11$ consecutive residuals vanish, so by Section 5 every later residual vanishes as well. $\square$

The entry claims the recurrence for $n > 3$. The theorem is at least as strong.

\section{Verification}

Four checks were run, each reproducible from the code accompanying this note, and the first two are what pin the reading of the entry's English.

\begin{itemize}

\item {\bfseries The published terms.} The walk count reproduces all 18 terms the entry publishes, exactly, in integer arithmetic, with the index taken from the entry's stated offset (0) rather than fitted. The first of them are 1, 11, 119, 1271, 13763, 149051, 1614359, 17485031,~\dots

\item {\bfseries The claim on the raw data.} Each claim was evaluated on the entry's published terms alone, with no model involved, wherever the proved range and the published data overlap. It holds there.

\item {\bfseries The annihilation test.} Carried to the derived bound $S = 11$ over the whole vector, in exact integer arithmetic, never sampled and never truncated.

\end{itemize}

\section{References}

\begin{enumerate}

\item OEIS Foundation Inc., \emph{The On-Line Encyclopedia of Integer Sequences}, entry \href{https://oeis.org/A125468}{A125468}, revision 17, last modified Aug 12 2026.

\item R. P. Stanley, \emph{Enumerative Combinatorics}, Vol.~1, 2nd ed., Cambridge University Press, 2012; \S4.7, the transfer-matrix method.

\item M. Kauers and P. Paule, \emph{The Concrete Tetrahedron}, Springer, 2011; C-finite sequences and their closure properties.

\item J. Berstel and C. Reutenauer, \emph{Noncommutative Rational Series with Applications}, Cambridge University Press, 2011; linear representations of counting automata and their reduction.

\end{enumerate}
\end{document}
